% arXiv submission — Frame Dependence: circuit robustness under task-preserving shift
% J. Melton, American Code Labs — September 2026
%
% DRAFT STATUS. Every number not yet measured is wrapped in \pending{} and renders
% as a red bracketed placeholder, so a partially-filled draft cannot be mistaken for
% a finished one at a glance. Before submission: \pendingcount must be zero and the
% \pendingnotice below must be deleted.
%
% The results and discussion sections are deliberately written WITHOUT committing to
% an outcome. The central measurement can land two ways and they support opposite
% prescriptions; prose that assumes one and gets edited later is how a paper ends up
% contradicting itself between sections.

\documentclass[11pt,a4paper]{article}

% ── Encoding & fonts ──────────────────────────────────────────────────────────
\usepackage[utf8]{inputenc}
\usepackage[T1]{fontenc}
\usepackage{lmodern}

% ── Math ──────────────────────────────────────────────────────────────────────
\usepackage{amsmath}
\usepackage{amssymb}

% ── Layout ────────────────────────────────────────────────────────────────────
\usepackage[margin=1.25in]{geometry}
\usepackage{microtype}
\usepackage{parskip}
\usepackage{setspace}

% ── Figures & tables ──────────────────────────────────────────────────────────
\usepackage{graphicx}
\usepackage{booktabs}
\usepackage{float}
\usepackage{caption}
\usepackage{array}
\usepackage{multirow}

% ── Code listings ─────────────────────────────────────────────────────────────
\usepackage{listings}
\usepackage{xcolor}

\definecolor{codegray}{rgb}{0.97,0.97,0.97}
\definecolor{commentgreen}{rgb}{0.0,0.5,0.0}
\definecolor{keywordblue}{rgb}{0.0,0.0,0.7}
\definecolor{stringbrown}{rgb}{0.5,0.1,0.0}

\lstdefinestyle{python}{
  language=Python, backgroundcolor=\color{codegray}, basicstyle=\small\ttfamily,
  keywordstyle=\color{keywordblue}\bfseries, commentstyle=\color{commentgreen},
  stringstyle=\color{stringbrown}, numberstyle=\tiny\color{gray}, numbers=left,
  numbersep=6pt, stepnumber=1, frame=single, framerule=0.4pt, xleftmargin=2em,
  breaklines=true, showstringspaces=false, tabsize=4, captionpos=b,
}

% ── Citations ─────────────────────────────────────────────────────────────────
\usepackage[numbers,sort&compress]{natbib}

% ── Hyperlinks ────────────────────────────────────────────────────────────────
\usepackage[colorlinks=true,linkcolor=blue,citecolor=blue,urlcolor=blue,
            pdftitle={Frame Dependence: Circuit Robustness Under Task-Preserving Distribution Shift},
            pdfauthor={J. Melton}]{hyperref}
\usepackage[htt]{hyphenat}
\emergencystretch=2em

% ── Draft placeholder machinery ───────────────────────────────────────────────
% A pending value is never rendered as a bare number. It is red, bracketed and
% counted, so that no draft can silently read as complete and no placeholder can be
% mistaken for a measurement in a screenshot or a skim.
\newcounter{pendingcount}
% Robust: a fragile macro inside \caption expands during .lof/.aux writing and
% crashes the engine. Declared robust pre-emptively -- captions are exactly where a
% pending value will eventually be needed.
\DeclareRobustCommand{\pending}[1]{\stepcounter{pendingcount}%
  \textcolor{red}{\textbf{[PENDING: #1]}}}

% The notice sits at the top of the document but must report a total that is only
% known at the bottom. Counting live would print 0 on the title page -- the exact
% "looks finished" signal this banner exists to prevent -- so the total is written to
% the .aux file at end of document and read back on the next pass.
\makeatletter
\newcommand{\pendingtotal}{\textbf{?}}
\AtEndDocument{\immediate\write\@auxout{%
  \string\gdef\string\pendingtotal{\thependingcount}}}
\makeatother

\newcommand{\pendingnotice}{%
  \begin{center}\fcolorbox{red}{yellow!15}{\parbox{0.92\textwidth}{\centering
  \textbf{\textcolor{red}{DRAFT --- \pendingtotal\ unmeasured values remain.}}\\
  \small Every red bracket marks a number the experiment has not yet returned.
  No value in this document is estimated, extrapolated or illustrative:
  a number is either measured or visibly absent.}}\end{center}\medskip}

\newcommand{\FD}{F(C_D \!\mid\! D)}
\newcommand{\FDp}{F(C_D \!\mid\! D')}

% ─────────────────────────────────────────────────────────────────────────────

\title{\textbf{Frame Dependence: Circuit Robustness\\Under Task-Preserving Distribution Shift}}

\author{
  J.\ Melton \\[4pt]
  American Code Labs \\
  \texttt{jmelton@americancode.org}
}

\date{September 2026 \\ \smallskip \small Preprint. Under review.}

% ─────────────────────────────────────────────────────────────────────────────
\begin{document}
\maketitle
\pendingnotice

% ── Abstract ──────────────────────────────────────────────────────────────────
\begin{abstract}
This paper measures whether a circuit discovered on a diverse set of sentence frames
remains faithful on frames it was not discovered on. Across \pending{n} task
families---indirect object identification, factual recall, subject--verb agreement and
in-context induction---on Llama-3.2-3B and Pythia-1.4B, we find \pending{headline:
does diverse discovery restore transfer, and does it depend on family}.
\pending{secondary: frame axis versus entity axis}.

The question is open because published results conflict. Merullo et al.\ report that
the GPT-2 Small indirect-object circuit ``generalizes surprisingly well'' to prompt
variants designed to break it. Rai et al.\ report cross-dataset faithfulness drops
reaching 79 percentage points under variations that preserve task semantics, and
conclude that circuit discovery recovers dataset-specific implementations. A third
study finds near-zero circuit overlap between template groups. One candidate
reconciliation is that the outcome is set by the \emph{discovery} data: a circuit found
on a narrow set has no opportunity to encode anything but that set---with a single
sentence frame, token position is perfectly confounded with syntactic role---while the
circuit Merullo et al.\ evaluate came from a fifteen-template protocol. If that is the
explanation, the remedy is a better discovery set rather than a new discovery
algorithm, which is what the field has so far proposed.

Deciding it requires a shift design that prior work does not supply: one holding the
number of templates fixed while changing only their identity. Comparing a one-frame
discovery set against a fifteen-frame evaluation set moves count and identity together
and cannot separate them. We discover a circuit on $D$ and score it on $D'$, where both
contain the same number of templates and share none, report the ratio $\FDp/\FD$ as
\emph{retention}, and admit a result only where the model still performs the task on
$D'$. Frame shift and entity shift are reported on separate axes rather than merged
into a single held-out number.

We release the protocol, twelve matched-diversity dataset pairs and reference
implementations, all runnable on a single consumer machine.
\end{abstract}

\tableofcontents
\newpage

% ─────────────────────────────────────────────────────────────────────────────
\section{Introduction}
\label{sec:intro}

Mechanistic interpretability is growing quickly and validating slowly. A survey of
indexed publications (\S\ref{sec:landscape}) finds \pending{N} papers matching
``mechanistic interpretability'', \pending{pct}\% of them from the last two years,
against \pending{N} matching ``circuit faithfulness'' and \pending{N} matching
``circuit generalization''. Circuits are being discovered far faster than they are
being tested, and the standard test---faithfulness---is normally computed on the data
the circuit was found on.

That is not a criticism of the metric. Faithfulness answers a well-posed question:
does this subgraph, on its own, reproduce the behaviour? The problem is what readers
take it to mean. A circuit reported as 87\% faithful is read as a claim about how the
model does the task, not about how the model does the task \emph{on these prompts},
and nothing in the usual reporting distinguishes those.

\subsection{Two Explanations With Opposite Consequences}

Our own prior work \citep{melton2026circuits} measured faithfulness falling from
$0.663$ to $0.228$ (Llama-3.2-3B) and $0.854$ to $-0.028$ (Pythia-1.4B) when circuits
discovered on one sentence frame were re-scored on fifteen, while both models
continued to perform the task. We reported the mechanism as unverified, and it is
worth being precise about why that caveat was unavoidable. The comparison moved two
variables together:

\begin{center}
\begin{tabular}{lll}
\toprule
 & discovery set & evaluation set \\
\midrule
templates & 1 & 15 \\
identity  & frame $f_1$ & frames $f_1 \dots f_{15}$ \\
\bottomrule
\end{tabular}
\end{center}

Two stories fit that evidence equally well.

\textbf{(a) Circuits are frame-specific.} A circuit found on any frame fails on other
frames. If so, a large and accelerating literature is reporting numbers that overstate
what has been established, and the remedy is procedural: faithfulness must be reported
on held-out data, always.

\textbf{(b) The discovery set was degenerate.} With a single frame, token position is
perfectly confounded with syntactic role. A head that appears circuit-critical may be
responding to a fixed position, and such a head fails anywhere the coincidence breaks.
If so, the published collapse is a fact about single-template discovery rather than
about circuits, and the remedy is upstream and cheap: discover on diverse data.

These prescribe opposite things---one says change how you report, the other says
change how you search---and no comparison that varies template count can choose
between them.

\subsection{Contributions}

Prior work has established that circuits often fail to transfer and that frame and
lexical shift are distinct axes; \S\ref{sec:positioning} states precisely which of our
earlier claims those results anticipate, and we do not restate them as new. What this
paper adds:

\begin{enumerate}
  \item \textbf{A test of the cheap remedy.} The field's response to dataset-specific
    circuits has been to propose new discovery algorithms. We test whether simply
    discovering on a diverse set is sufficient, which is the intervention a
    practitioner would try first and which no prior study isolates
    (\S\ref{sec:protocol}).
  \item \textbf{Matched-diversity disjoint splits.} Template count held fixed, template
    identity disjoint---the design that separates ``the frames differ'' from ``there
    are more frames''. Prior variation axes are systematic transformations or lexical
    substitutions and do not make this separation (\S\ref{sec:axes}).
  \item \textbf{Retention as the reported quantity.} $\FDp/\FD$ rather than $\FDp$,
    which distinguishes a circuit that was weak everywhere from one that collapsed off
    its discovery set (\S\ref{sec:protocol}).
  \item \textbf{Measurements} across \pending{n} families and two models, one of them
    a family not covered at this scale by prior work (\S\ref{sec:results}).
  \item \textbf{Two harness defects, disclosed} (\S\ref{sec:defects}), both silent on
    near-uniform-length data and both capable of contaminating this comparison.
\end{enumerate}

% ─────────────────────────────────────────────────────────────────────────────
\section{Background}
\label{sec:background}

\subsection{Circuits and Faithfulness}

Following \citet{wang2022ioi}, a circuit $C$ is a set of model components---here
attention heads---hypothesised to implement a behaviour. Faithfulness mean-ablates
every component outside $C$ and measures the recovered logit difference, normalised
against the fully-ablated floor:
\[
  F(C) \;=\; \frac{\mathrm{LD}(\text{only } C \text{ active}) - \mathrm{LD}_{\text{floor}}}
                  {\mathrm{LD}(\text{clean}) - \mathrm{LD}_{\text{floor}}}
\]
The floor term matters. An unnormalised ratio implicitly treats the fully-ablated
model as scoring zero, which is not generally true and can be negative.

\subsection{Field Context}
\label{sec:landscape}

Validation of circuits has not kept pace with their discovery. Indexed publication
counts (OpenAlex, exact-phrase title-and-abstract search, retrieved 2026-08-13) place
``mechanistic interpretability'' at 2{,}957 works, 87\% of them from 2025--2026, and
``activation patching'' at 561. Terms naming the validation question are two orders of
magnitude rarer.

\begin{table}[H]
  \centering
  \caption{Exact-phrase publication counts, OpenAlex, 2026-08-13. Counts use quoted
    phrase search; the default fuzzy search is not phrase search and inflates
    \texttt{activation patching} from 561 to 41{,}006. Terms whose usage is dominated
    by electrical-circuit senses (\texttt{circuit validation}, \texttt{circuit
    robustness}) are excluded.}
  \label{tab:landscape}
  \smallskip
  \begin{tabular}{lrrr}
    \toprule
    phrase & all & 2025--26 & recent \\
    \midrule
    mechanistic interpretability  & 2{,}957 & 2{,}578 & 87\% \\
    activation patching           &     561 &     442 & 79\% \\
    circuit generalization        &       8 &       1 & 12\% \\
    circuit faithfulness          &       8 &       5 & 62\% \\
    \bottomrule
  \end{tabular}
\end{table}

These counts measure vocabulary rather than activity and should be read only as an
order of magnitude. The three studies most directly relevant to this paper
(\S\ref{sec:related}) use none of the phrases above, describing the same problem as
adaptive circuit behaviour, dataset-specific circuits and rephrasing variance. The
literature review in \S\ref{sec:related}, not Table~\ref{tab:landscape}, establishes
what is known.

\subsection{Related Work: The Field Disagrees With Itself}
\label{sec:related}

Circuit robustness is not an unexamined question. It is an examined question with
contradictory answers, and that contradiction is what this paper is positioned against.

\textbf{Evidence that circuits generalize.} \citet{merullo2024adaptive} test the GPT-2
Small IOI circuit on prompt variants that violate the assumptions of the published
algorithm, and report that it ``generalizes surprisingly well, reusing all of its
components and mechanisms while only adding additional input edges''. They identify a
previously undescribed mechanism, S2 Hacking, that accounts for the circuit succeeding
on formats where the original algorithm should fail, and conclude that circuits may be
more flexible than previously recognised.

\textbf{Evidence that circuits do not generalize.} \citet{rai2026dcd} test the
single-circuit-per-task assumption across four tasks (indirect object identification,
entity binding, arithmetic addition, sequence completion) on GPT-2 Small,
Qwen2.5-7B-Instruct and Llama-3.1-8B-Instruct, varying datasets along complexity,
syntax and domain axes while preserving task semantics. They report in-distribution to
cross-dataset faithfulness drops reaching 79 percentage points at circuit size 0.05,
and conclude that existing methods recover dataset-specific rather than task-general
circuits. Independently, \citet{anon2026variance} study rephrasing variance and find
that different templates activate substantially different circuits, with near-zero
overlap between template groups, across GPT-2 and Pythia models on subject--verb
agreement, IOI, greater-than and coding tasks.

\textbf{These cannot all be straightforwardly true.} One candidate reconciliation is
that the answer depends on the discovery set. \citet{merullo2024adaptive} evaluate a
circuit produced by the Wang et al.\ protocol, which uses fifteen templates with
balanced name ordering; a circuit discovered on a narrow set has no opportunity to
encode anything but the narrow set. If that is the reconciliation, frame dependence is
a property of \emph{how the circuit was found} rather than of circuits as such, and
the remedy is cheaper than it currently appears.

\section{The Protocol}
\label{sec:protocol}

\subsection{Definitions}

Let $\mathcal{T}$ be a task, and let $D$ and $D'$ be two datasets that both instantiate
$\mathcal{T}$. Discover $C_D$ on $D$ by any procedure; here, the combined
patching-and-ablation ranking of \S\ref{sec:discovery}. Then report:
\[
  \text{retention } R \;=\; \frac{\FDp}{\FD}
  \qquad\qquad
  \text{frame dependence } \Delta \;=\; \FD - \FDp
\]
$R$ is the headline because it is scale-free; $\Delta$ is reported alongside because
$R$ is unstable when $\FD$ is near zero.

\subsection{Matched Diversity}

$D$ and $D'$ must have the same number of templates and share none. This is the whole
design. Holding count fixed removes the confound of \S\ref{sec:intro}; sharing no
frames is what makes the shift a shift.

We additionally hold clause \emph{shape} constant across the two template banks where
the family permits it. Varying syntax as well would broaden the shift, but it would
also reintroduce a second moving part, and the point of this protocol is that exactly
one thing moves.

\subsection{Shift Axes}
\label{sec:axes}

\begin{center}
\begin{tabular}{lll}
\toprule
axis & templates & entity pool \\
\midrule
frame  & disjoint & shared \\
entity & shared   & disjoint \\
both   & disjoint & disjoint \\
\bottomrule
\end{tabular}
\end{center}

Reporting only a merged ``held-out'' number would conflate two different objects. A
circuit that survives new names but not new frames is a syntactic-position circuit; a
circuit that survives new frames but not new names has memorised lexical items. Both
are interesting; they are not the same finding.

\subsection{The Task-Preservation Gate}
\label{sec:gate}

A faithfulness collapse on $D'$ means nothing about circuits if the model cannot do
the task on $D'$. We therefore report clean logit difference on both splits and admit
a row only if behaviour is retained above a threshold of $0.8$; failing rows are
listed as excluded, with their numbers, rather than dropped silently or averaged in.

% ─────────────────────────────────────────────────────────────────────────────
\section{Datasets}
\label{sec:datasets}

Four families, each in three shift axes, each split into $D$ and $D'$ of 200 examples:

\begin{center}
\begin{tabular}{llll}
\toprule
family & mechanism exercised & templates/bank & answer contrast \\
\midrule
\texttt{ioi}       & routing / name suppression      & 15 & indirect object vs subject \\
\texttt{factual}   & retrieval                       & 6  & capital vs other capital \\
\texttt{agreement} & long-range syntactic dependency & 6  & \emph{are} vs \emph{is} \\
\texttt{induction} & in-context copying              & 6  & paired vs unpaired token \\
\bottomrule
\end{tabular}
\end{center}

Bank A for \texttt{ioi} reproduces the fifteen frames of \citet{melton2026circuits}
exactly, so $D$ matches the previously published evaluation set and the new $D'$ is
the only novel object.

\subsection{Construction Constraints}

Three constraints are enforced at generation time; each was added because violating it
produces a plausible-looking but wrong measurement rather than an error.

\textbf{Single-token answers.} The metric reads one token id at the final position, so
a multi-token answer would be scored on its first piece---often a shared prefix, and
therefore not diagnostic. Candidate answers are filtered to single tokens under
\emph{both} tokenizers.

\textbf{Clean/corrupt token-length parity.} Activation patching caches per-head output
on a clean pass and splices it into a corrupt pass at the same tensor positions. That
correspondence requires position $i$ to mean the same thing in both. IOI corruption
swaps two single-token names and is length-preserving by construction; entity
substitution in other families is not. Pairs failing parity under either tokenizer are
dropped---\pending{pct}\% of \texttt{factual} candidates, and none elsewhere.

\textbf{Template-balanced sampling.} The parity filter does not drop uniformly across
frames, so taking the first $n$ survivors would skew the template mix and confound
frame identity with frame frequency. Records are taken round-robin across template ids.

\textbf{Entity pools are filtered before splitting, not after.} Splitting into fixed
banks and filtering each independently lets tokenizer attrition fall unevenly; in an
early version this left one induction bank with five usable words and the other with
one. An unequal pool is itself a difference between $D$ and $D'$.

% ─────────────────────────────────────────────────────────────────────────────
\section{Method}
\label{sec:method}

\subsection{Discovery}
\label{sec:discovery}

Two per-head metrics over a full $\text{layers} \times \text{heads}$ sweep:

\emph{Patching.} Run the corrupt prompt with head $(l,h)$'s clean output spliced in,
and measure the fraction of the clean-corrupt gap recovered. High means the head
carries task signal.

\emph{Ablation.} Mean-ablate $(l,h)$ on the clean prompt and measure the logit-difference
drop. High means nothing else compensates for its loss.

Each is min-max normalised across all heads and summed; the circuit is the top ten by
that sum. The two metrics disagree often enough to matter, which is why the selection
uses both.

\subsection{Implementation Requirements}
\label{sec:defects}

Two implementation properties are required for this comparison and are not guaranteed
by a harness built for a single near-uniform dataset. Both are silent when violated:
they produce numbers rather than errors.

\textbf{The mean-ablation reference must be computed over the full example set.} A
caching step that assigns rather than accumulates leaves only the final batch's mean,
so the ablation value is estimated from a fraction of the data. The quantity estimated
is the same, but its variance is not, and $D$ and $D'$ must be scored against the same
reference for their difference to mean anything.

\textbf{Sequences must be padded to one global length, not per batch.} The cached mean
carries a sequence dimension, so a mean cached at one padded length cannot be broadcast
onto a batch padded to another. Datasets of near-uniform length hide this; three of the
four families here vary enough in length to expose it.

The harness used for all measurements below satisfies both. Figures produced under
these corrections do not reproduce previously published IOI values to the last decimal,
and the two are not presented in a common column.

\section{Results}
\label{sec:results}

\begin{table}[H]
  \centering
  \caption{Retention under frame shift. $R = \FDp/\FD$. ``behav.'' is clean logit
    difference on $D'$ as a fraction of $D$; rows below $0.8$ fail the
    task-preservation gate and are excluded from any aggregate.}
  \label{tab:frame}
  \smallskip
  \begin{tabular}{llccccc}
    \toprule
    model & family & $\FD$ & $\FDp$ & $R$ & behav. & gate \\
    \midrule
    Llama-3.2-3B & \texttt{ioi}       & \pending{} & \pending{} & \pending{} & \pending{} & \pending{} \\
    Llama-3.2-3B & \texttt{factual}   & \pending{} & \pending{} & \pending{} & \pending{} & \pending{} \\
    Llama-3.2-3B & \texttt{agreement} & \pending{} & \pending{} & \pending{} & \pending{} & \pending{} \\
    Llama-3.2-3B & \texttt{induction} & \pending{} & \pending{} & \pending{} & \pending{} & \pending{} \\
    \midrule
    Pythia-1.4B  & \texttt{ioi}       & \pending{} & \pending{} & \pending{} & \pending{} & \pending{} \\
    Pythia-1.4B  & \texttt{factual}   & \pending{} & \pending{} & \pending{} & \pending{} & \pending{} \\
    Pythia-1.4B  & \texttt{agreement} & \pending{} & \pending{} & \pending{} & \pending{} & \pending{} \\
    Pythia-1.4B  & \texttt{induction} & \pending{} & \pending{} & \pending{} & \pending{} & \pending{} \\
    \bottomrule
  \end{tabular}
\end{table}

\begin{table}[H]
  \centering
  \caption{Frame versus entity shift, IOI only.}
  \label{tab:axes}
  \smallskip
  \begin{tabular}{lcccc}
    \toprule
    model & $R$ (frame) & $R$ (entity) & $R$ (both) & $\FD$ \\
    \midrule
    Llama-3.2-3B & \pending{} & \pending{} & \pending{} & \pending{} \\
    Pythia-1.4B  & \pending{} & \pending{} & \pending{} & \pending{} \\
    \bottomrule
  \end{tabular}
\end{table}

\subsection{Circuit Overlap}

\pending{How much do $C_D$ and a circuit discovered directly on $D'$ overlap? A
circuit that fails to transfer but whose replacement shares most of its heads is a
different phenomenon from one replaced wholesale.}

% ─────────────────────────────────────────────────────────────────────────────
\section{Analysis}
\label{sec:analysis}

\emph{Drafting note --- to be written once Table~\ref{tab:frame} is filled, and to be
written once. The measurement selects which of these the paper argues; it is not a
menu to hedge across. Both branches share the framing above, so no section written so
far needs revisiting either way.}

\textbf{If retention is high} ($R$ near $1$ across families): the published collapse
was about single-template discovery, not about circuits. The prescription is upstream
and cheap, the finding bounds our own prior paper, and the protocol's value is as a
cheap check that a discovery set was not degenerate.

\textbf{If retention is low}: frame dependence is real, general across mechanisms as
different as retrieval and induction, and faithfulness reported in-distribution
systematically overstates what has been shown.

\textbf{If retention splits by family}: the most informative outcome, and the one that
would need the most careful writing---which mechanisms are portable and which are
frame-bound is a stronger claim than either uniform result.

% ─────────────────────────────────────────────────────────────────────────────
\section{Limitations}
\label{sec:limitations}

\textbf{Two models, both small.} Llama-3.2-3B and Pythia-1.4B. Whether frame
dependence changes with scale is untested here.

\textbf{Attention heads only.} MLP components are not in the circuit vocabulary, so a
circuit that appears to lose faithfulness may have had its work taken over by
components this method cannot see.

\textbf{One discovery procedure.} Retention is a property of the circuit-finding
method as much as of the model. A different procedure may produce more portable
circuits, and this protocol is the way to check that rather than an argument that it
would not.

\textbf{Templates are hand-written.} Frame banks were authored, not sampled from a
corpus, so ``frame shift'' means the shift the author constructed.

\textbf{The shift is synthetic.} $D'$ is task-preserving by construction and gated on
behaviour, but it is not natural text.

% ─────────────────────────────────────────────────────────────────────────────
\section{Reproducibility}
\label{sec:repro}

Code and data: \url{https://github.com/american-code/ResearchPapers}. Archived at
\pending{Zenodo concept DOI}.

\begin{center}
\begin{tabular}{ll}
\toprule
artifact & path \\
\midrule
dataset generator      & \texttt{data/circuit-robustness/generate\_shift\_pairs.py} \\
intervention library   & \texttt{data/circuit-robustness/circuit\_lib.py} \\
discovery              & \texttt{data/circuit-robustness/discover\_circuit.py} \\
scoring                & \texttt{data/circuit-robustness/score\_faithfulness.py} \\
end-to-end protocol    & \texttt{data/circuit-robustness/run\_protocol.sh} \\
generated splits       & \texttt{data/circuit-robustness/*-\{D,Dprime\}.json} \\
results                & \texttt{data/circuit-robustness/results/} \\
\bottomrule
\end{tabular}
\end{center}

All runs are seeded (42) and executed on a single Apple M1 Max with 32\,GB unified
memory via MLX.

% ─────────────────────────────────────────────────────────────────────────────
\bibliographystyle{plainnat}
\bibliography{refs}

\end{document}
